% Exposé XIII. Conclusion (written for the paper; scope of every claim as in theory/propositions.json, public_statement).
\section{Conclusion}\label{sec:conclusion}

Fine-tuning moves a point. Treating a PEFT method as a pointed map $\rho:(Q,q_0)\to(\Theta,\theta_0)$ into weight space, and studying the map more than the module that implements it, turns many informal comparisons between methods into questions with exact answers. Three invariants of the map do most of the work. The image germ at $\theta_0$ governs what a method can reach, through its first-order image $S_1$, its expressive dimension $d(M)$, its cut-rank and the type of its base point. The fibres, together with the optimizer's metric, govern how it trains. Base change governs how it travels when it is initialised by splitting, quantised, restarted or merged. Extensions such as adapters and prefixes add one further layer, where the network itself enters and merging becomes a lifting problem.

Several results follow, each with the scope of its proof. Zero-initialised LoRA starts at the apex of a cone, and so does HRA at the paired initialisation of HF PEFT when $W_0$ has full column rank and $r\le n-2$ is even (\thmref{Theorem}{II.2}, \thmref{Theorem}{II.7}). No exact initialisation lets LoRA reach an update of rank above $2r$ (\thmref{Theorem}{II.1}). Exactly orthogonal methods never change the singular values of a weight in exact arithmetic, however often they are merged (\thmref{Corollary}{II.6}, \thmref{Theorem}{V.3}). Rank is a cut and dimension is a volume, so at equal budget a structured method can gain rank at the same dimension, as GraLoRA does, or trade dimension for rank, as LoHa does when $2r^2<m+n-1$ and $r\ge3$ (\thmref{Proposition}{II.12}). Under Adam with no $\varepsilon$, weight decay or clipping and one shared schedule, LoRA+ follows exactly the trajectory of LoRA with $\alpha$ multiplied by the learning-rate ratio (\thmref{Corollary}{III.11}; the LoRA/Adam case is due to \citealp{schulman2025lora}), and on LoRA's full-rank locus Adam with channel-uniform hyperparameters respects only the signed permutations inside the $\GL_r$ gauge (\thmref{Theorem}{III.12}). Whether a scaling adapter folds into the weights before it depends on the activation, and a local rewrite calculus certifies when it does (\thmref{Theorem}{VI.3}). The atlas reads the same coordinates off \nMethods{} methods from \nPapers{} papers and records \nArrows{} simulation arrows and \nObstructions{} obstructions, each backed by an explicit map or a named test.

Several of these results are short once the setting is in place, which is what the rising sea on the title page refers to (Appendix~\ref{app:making-of}). In the slice, re-pointing an additive method sweeps out the difference set $C-C$ of its update shape, and for LoRA$_r$ that set is the matrices of rank at most $2r$. The bound on every exact initialisation is then a few lines (\thmref{Theorem}{II.1}). Once LoRA's factors are seen to carry a torus of rescalings that fixes the weight trajectory, LoRA+ is a point on that torus (\thmref{Corollary}{III.11}). In the lens, GaLore without weight decay becomes one-sided LoRA on a moving base by a one-line induction (\thmref{Proposition}{IV.2}). Read off a tensor diagram, the maximum rank of a structured adapter is a min-cut (\thmref{Theorem}{II.10}). Three of these four facts were known before this paper, from the PiSSA-to-LoRA conversion \citep{meng2024pissa}, \citet{schulman2025lora} and \citet{xtorrobahennigen2025duality}. What the setting adds is a short proof, an exact scope and the same proof for every method of the same shape. Other results did not dissolve. The classification of initialisations (\thmref{Theorem}{II.2}) and rebasing closure (\thmref{Theorem}{V.3}) required real algebraic geometry and Lie theory, and their proofs say where.

For practitioners these statements become checks that can be made before training. Compare methods at equal expressive dimension and cut-rank as well as at equal budget. Ask whether an initialisation is an apex, and which part of its frame the optimizer will see (\thmref{Corollary}{III.13}). Before treating a new hyperparameter as a new method, check whether it lies on a hyperparameter torus. Before merging adapters factor by factor, check that the rule is constant on factorisation fibres (\thmref{Proposition}{V.4}). Before relying on an orthogonal method to keep a spectrum, check that its implementation is exactly orthogonal. Implementations matter throughout. Several results distinguish a method as published from its HF PEFT implementation, such as OFT's Cayley--Neumann default, HRA's paired initialisation and LoftQ's treatment of the scale $\alpha/r$ (\thmref{Proposition}{V.2}), and the catalogue records which form each statement covers.

The limits are part of the statements. Expressivity is not learnability, and only Exposé~\ref{sec:fibres} is dynamical. Its conservation laws hold under gradient flow, drift under discrete steps by an amount of second order in the step size relative to the charge, and do not cover Adam or momentum. The classification of initialisations is set-level and in exact arithmetic. Mergeability is certified box by box under explicit side conditions, and the rewrite rules give sufficient conditions and local obstructions; they are not a decision procedure for whole networks. Words such as Noether, Erlangen, gauge and moduli are analogies beyond what is proved, and Exposé~\ref{sec:analogy} says where each one stops.

The falsifiable predictions of Exposé~\ref{sec:predictions} say where the theory can be tested, each with its optimizer scope. The open problems of Exposé~\ref{sec:analogy} say where it should grow next: a dynamical classification of initialisations modulo the optimizer's shadow of the gauge, a conserved or monotone quantity for sign-like updates such as Adam's, the frontier of pairwise incomparable methods at a fixed budget, and whole-network statements about prefix tuning. We hope the frame serves as a habit of asking three questions of any new method: what its map can reach, how its fibres meet the optimizer, and what happens to it when the base point moves.
